\chapter{Будущее: клеточные комплексы как язык для составных актов}
\label{ch:cellcomplex}

\section{Зачем вообще думать о клетках}

Гиперграфовая ветвь уже показывает, что попарных рёбер иногда недостаточно. Но и
гиперребро — ещё не вся картина. У памяти есть не только набор совместно действующих
посылок, но и \emph{структура их совместности}: границы, циклы, составные акты,
которые хочется видеть не как мешок из входов, а как объект более высокой размерности.

\section{Базовая идея}

Клеточный комплекс позволяет держать объекты разных размерностей:
\begin{enumerate}
\item 0-клетки как узлы памяти;
\item 1-клетки как обычные рёбра;
\item 2-клетки как составные акты, выводы или сцены;
\item при желании и более высокие клетки для сложных конфигураций.
\end{enumerate}

Это полезно тогда, когда память должна различать не только ``что с чем связано'',
но и ``что из чего составлено как единый акт''.

\section{Что это даёт памяти}

\paragraph{Во-первых,} можно хранить составную ситуацию как один объект, не
раскладывая её полностью на попарные связи.

\paragraph{Во-вторых,} появляются естественные границы: если 2-клетка кодирует
составной вывод или эпизодический акт, её boundary прямо указывает на набор посылок
или компонент.

\paragraph{В-третьих,} становится богаче язык локальной ревизии: можно понижать не
только узел или ребро, но и целую составную конфигурацию.

\section{Чего это не решает}

\begin{pitfall}[Новая структура не заменяет дисциплину provenance]
Сам по себе переход к клеткам не делает память ни истиннее, ни объяснимее. Если у
системы нет хорошего provenance, gates, deletion-path и ревизии, клеточный комплекс
лишь переносит старые проблемы в более красивую геометрию.
\end{pitfall}

\section{Правильный статус этой ветви}

Для текущего проекта клеточные комплексы — это язык \emph{будущего структурного
обобщения}, а не ближайшая реализация. Их стоит читать так:
\begin{enumerate}
\item как концептуальный мост между графом и n-арными актами;
\item как способ точнее думать о границах составных выводов;
\item как материал для дальнего сравнения с гиперграфами и динамической памятью.
\end{enumerate}

Именно в таком статусе эта ветвь совместима с нынешним roadmap: не как срочная
замена субстрата, а как более богатая геометрия для будущего поколения памяти.
